¿Qué son las herramientas forenses? ¿Para qué sirven? Da un ejemplo y explica su uso.

Son herramientas de software o hardware especializadas en adquirir, preservar y analizar evidencias digitales procedentes de computadoras, redes, sistemas operativos y memoria tras un incidente de seguridad o delito informático, garantizando que las evidencias no se vean alteradas.

Las herramientas forenses sirven para reconstruir los eventos cibernéticos de un incidente para que el especialista en ciberseguridad comprenda cómo se vulneró el sistema, cuándo se vulneró y qué información se vio comprometida durante el ataque, de manera que se puedan recuperar evidencias legalmente válidas y admisibles ante un tribunal, así como reducir el daño producido por un ataque.

Un ejemplo de herramienta forense es FTK (\textit{Forensic Toolkit}), la cual está orientada al procesamiento y análisis de volúmenes de evidencia digital. FTK se utiliza para examinar discos, imágenes forenses, archivos, correos electrónicos y datos recopilados durante una investigación. Además, incorpora funciones para trabajar con evidencia cifrada y credenciales, así como para verificar la integridad de los datos manteniendo la trazabilidad de las evidencias desde la adquisición hasta la elaboración del informe forense~\cite{HF}.
